\documentclass[../../../include/open-logic-section]{subfiles}

\begin{document}

\olfileid{sth}{replacement}{ref}
\olsection{تعويض او انعکاس}

د \stagesinex له لارې د تعويض د توجيې زموږ وروستۍ هڅه له يوې ژورې
او ښکلې پايلې پيلېږي:\footnote{يادونه: ټول فورمولونه پاراميټرونه
لرلے شي، مګر چې په څرګند ډول بل څه ويل شوي وي.}
%In this, I will use overlining, such as $x_1, \ldots, x_n$, to
%abbreviate ``$x_1, \ldots, x_n$'': 
\begin{thm}[د انعکاس شېما]\ollabel{reflectionschema}
د هر فورمول $\phi$ دپاره:
\[
\forall \alpha \exists \beta > \alpha (\forall x_1 \ldots, x_n \in
V_\beta)(\phi(x_1, \ldots, x_n) \liff \phi^{V_\beta}(x_1, \ldots, x_n))
\]
\end{thm}
\noindent 
لکه په \olref[sth][replacement][strength]{formularelativization} کښې،
$\phi^{V_\beta}$ د $\phi$ د ټولو کميت ټاکونکو د $V_\beta$ سټ ته
د محدودولو پايله ده۔ نو انعکاس په بداهتي ډول داسې وايي: که $\phi$
په ټوله سلسله کښې رښتيا وي، نو $\phi$ د سلسلې په هر څومره لرې
\emph{پيلنيو برخو} کښې هم رښتيا وي۔

\citet{Montague1961} او \citet{Levy1960} وښودله چې د تعويض او
انعکاس مناسبې رسمي بڼې د~$\Z$ په نسبت همارزښته دي؛ نو له دغو دوو
هر يو چې ورزيات کړو، $\ZF$ ترلاسه کېږي۔ (دا پايلې په
\olref[sth][replacement][refproofs]{sec} کښې ثابتوو۔) له دې
همارزښت څخه هيله کېدے شي چې د \stagesinex{} په مټ انعکاس او تعويض
داسې توجيه کړو: د \stagesinex په منلو سره سلسله بايد ډېره، ډېره
لوړه وي؛ دومره لوړه چې د هغې په اړه هېڅ وينا ئې د لوړوالي د حد
ټاکلو دپاره بس نۀ وي۔ دا داسې هم پوهېدے شو چې که کومه جمله~$\phi$
په ټوله سلسله کښې رښتيا وي، نو د سلسلې په هر څومره لرې پيلنيو
برخو کښې هم رښتيا وي۔ همدا انعکاس دے۔

بيا هم، دا د تعويض د دننۍ توجيې يوه رښتيا هيله‌بښونکې هڅه ښکاري۔
خو دلته د دې په اړه تر دې ډېر څه ويل پکار دي۔ اوس پخپله پرېکړه
وکړئ چې دا هڅه بريالۍ ده که نۀ۔\footnote{کېدے شي د نورو معلوماتو
دپاره \citet[95--100]{Incurvati2020} هم ولولئ.}

\end{document}
